\chapter*{Daftar Simbol}


\begin{center}
\small % Kurangi ukuran font
\renewcommand{\arraystretch}{1.05} % Kurangi jarak baris dari 1.2
\setlength{\tabcolsep}{4pt} % Kurangi padding dari 6pt
\begin{tabular}{p{2.5cm}p{0.2cm}p{10cm}} % Kurangi lebar kolom pertama
$G = (V, E)$ & : & Graf berarah yang merepresentasikan ekonomi gim dengan himpunan vertex $V$ dan himpunan edge $E$ \\
$V$ & : & Himpunan vertex (node) dalam graf ekonomi gim \\
$E$ & : & Himpunan edge (sisi) berarah dalam graf \\
$v_i$ & : & Vertex ke-$i$ dalam graf, merepresentasikan komponen ekonomi \\
$n$ & : & Jumlah node dalam graf ekonomi gim \\
$w_{ij}$ & : & Bobot aliran sumber daya dari node $i$ ke node $j$ \\
$W$ & : & Matriks bobot aliran antar node: $W = [w_{ij}]_{n \times n}$ \\
$\text{Outflow}(v_i)$ & : & Jumlah total aliran keluar dari node $v_i$: $\sum_{j=1}^{n} w_{ij}$ \\
$\text{Inflow}(v_i)$ & : & Jumlah total aliran masuk ke node $v_i$: $\sum_{j=1}^{n} w_{ji}$ \\
$f(G)$ & : & Fungsi objektif untuk optimasi keseimbangan graf \\
$P_0$ & : & Populasi awal dalam algoritma evolusioner \\
$P_t$ & : & Populasi pada generasi ke-$t$ \\
$P_{t+1}$ & : & Populasi pada generasi ke-$(t+1)$ hasil evolusi \\
$x_i$ & : & Individu (solusi kandidat) ke-$i$ dalam populasi \\
$x^*$ & : & Solusi optimal hasil algoritma evolusioner \\
$f_{\text{gen}}(G)$ & : & Fungsi fitness pada tahap \textit{generator} untuk mengevaluasi validitas struktur graf ekonomi gim \\
$\mathrm{validate}(v)$ & : & Fungsi bantu yang menghitung jumlah batasan yang tidak terpenuhi oleh node $v$ \\
$f_1(s^{p_j}_t, x)$ & : & Fungsi fitness pada tahap \textit{balancer} untuk mengukur kedekatan nilai sumber daya terhadap target absolut \\
$s^{p_j}_t$ & : & Jumlah sumber daya dalam pool $p_j$ pada waktu ke-$t$ hasil simulasi ekonomi gim \\
$p_j$ & : & Pool sumber daya ke-$j$ dalam graf ekonomi gim \\
$x$ & : & Nilai target absolut yang ingin dicapai oleh sumber daya (misal jumlah \textit{gold}) \\
$\alpha$ & : & Parameter ambang batas (\textit{threshold}) untuk toleransi rata-rata hasil simulasi pada fungsi fitness \\
$m$ & : & Jumlah simulasi yang dijalankan untuk mengurangi pengaruh keacakan (\textit{randomness}) \\
$i$ & : & Indeks simulasi ke-$i$ dalam perhitungan rata-rata fitness \\
$\mathrm{prop}(s^{p_j}_t, x)$ & : & Fungsi proporsional untuk mengukur kedekatan nilai sumber daya terhadap target \\
$N$ & : & Ukuran populasi dalam algoritma evolusioner \\
$f(\cdot)$ & : & Fungsi fitness untuk evaluasi kualitas individu \\
$p_c$ & : & Probabilitas crossover dalam algoritma evolusioner \\
$p_m$ & : & Probabilitas mutasi dalam algoritma evolusioner \\
$x_i'$ & : & Individu hasil mutasi dari $x_i$ \\
$\epsilon$ & : & Variabel acak dari distribusi normal untuk mutasi \\
$\sigma^2$ & : & Variansi untuk operasi mutasi \\
$\mathcal{N}(0, \sigma^2)$ & : & Distribusi normal dengan mean 0 dan variansi $\sigma^2$ \\
\end{tabular}
\end{center}